% SPDX-FileCopyrightText: 2026 Will Cook
% SPDX-License-Identifier: CC-BY-4.0
%
% One of the Erdős Problem Notes: a short standalone treatment of a single
% problem in the ErdosProblems expansion library.  Build with
% `tectonic erdos-251-prime-gap-dyadic-series.tex` or `make`.
\documentclass[11pt]{article}
% SPDX-FileCopyrightText: 2026 Will Cook
% SPDX-License-Identifier: CC-BY-4.0
%
% Shared preamble for the Erdős Problem Notes: the short, standalone,
% problem-owned manuscripts covering the expansion library `ErdosProblems`.
%
% One note owns one Erdős problem.  Each note repeats whatever context its own
% reader needs, so that a reader who arrives at a single problem never has to
% assemble it from the gateway paper.  Deliberate repetition across notes is the
% design, not an oversight.
%
% Source links resolve against one pinned formal-source commit, declared once
% below.  `scripts/check_problem_note_sources.py` verifies that every authored
% (file, line, declaration) triple names that exact declaration in the pinned
% snapshot, so a link cannot silently rot as later work moves lines.

\usepackage{paper-house-style}
\usepackage{array}
\usepackage{booktabs}
\usepackage{enumitem}
\setlist{topsep=0.35em,itemsep=0.2em}
\usepackage[colorlinks=true,linkcolor=linkinternal,urlcolor=linkexternal,
  citecolor=linkinternal]{hyperref}
\hypersetup{bookmarksnumbered=true,bookmarksopen=true,bookmarksopenlevel=1,
  pdfauthor={Will Cook},pdflang={en-GB}}

% ---------------------------------------------------------------------------
% Pinned formal source.  \commit names the checkpoint every link resolves
% against; the checker reads that snapshot rather than the working tree, so the
% printed coordinates stay correct however later waves move declarations.
% ---------------------------------------------------------------------------
\newcommand{\commit}{e5fd26a6d1b1a346430205eab5385b4c9565d004}
\newcommand{\commitshort}{e5fd26a6d1b1}
\newcommand{\repobase}{https://github.com/wcook04/plectis-lean-erdos249-257/blob/\commit}
\newcommand{\sourceurl}{https://github.com/wcook04/plectis-lean-erdos249-257/tree/\commit}
\newcommand{\PX}{ErdosProblems}

% Declaration names stay inside link targets.  The printed label is either a
% spaced, lower-case reading of the declaration name (\lref) or a short authored
% phrase (\lword); neither prints a module path or a raw Lean identifier.
\ExplSyntaxOn
\cs_new_protected:Npn \noteleanlabel #1
  {
    \tl_set:Nx \l_tmpa_tl { \tl_to_str:n {#1} }
    \regex_replace_all:nnN { _+ } { \c{space} } \l_tmpa_tl
    \regex_replace_all:nnN
      { ([a-z0-9])([A-Z]) }
      { \1\c{space}\2 }
      \l_tmpa_tl
    \tl_set:Nx \l_tmpa_tl { \tl_lower_case:n { \tl_use:N \l_tmpa_tl } }
    \tl_use:N \l_tmpa_tl
  }
\ExplSyntaxOff
\newcommand{\leanlabel}[1]{\noteleanlabel{#1}}
\newcommand{\lref}[3]{\href{\repobase/\PX/#1\#L#2}{\leanlabel{#3}}}
\newcommand{\lword}[4]{\href{\repobase/\PX/#1\#L#2}{#4}}
\newcommand{\lloc}[2]{\href{\repobase/\PX/#1\#L#2}{Lean source}}

% Links into a sibling library, named repository-relative.  The expansion
% library is implicit in \lref/\lword, because that is where problem-owned
% work lands.  Several problems have their headline theorem in the older
% reviewed #249/#257 corpus, and a note that could not name that library
% would have to paraphrase a checked statement instead of linking it.
\newcommand{\mref}[3]{\href{\repobase/#1\#L#2}{\leanlabel{#3}}}
\newcommand{\mword}[4]{\href{\repobase/#1\#L#2}{#4}}
\newcommand{\mloc}[2]{\href{\repobase/#1\#L#2}{Lean source}}
\emergencystretch=2em

\newtheorem{theorem}{Theorem}[section]
\newtheorem{proposition}[theorem]{Proposition}
\newtheorem{lemma}[theorem]{Lemma}
\newtheorem{corollary}[theorem]{Corollary}
\theoremstyle{definition}
\newtheorem{definition}[theorem]{Definition}
\newtheorem{example}[theorem]{Example}
\newtheorem{problem}[theorem]{Problem}
\theoremstyle{remark}
\newtheorem*{remark}{Remark}

\newcommand{\N}{\mathbb{N}}
\newcommand{\Npos}{\mathbb{N}_{>0}}
\newcommand{\Z}{\mathbb{Z}}
\newcommand{\Q}{\mathbb{Q}}
\newcommand{\R}{\mathbb{R}}
\newcommand{\lcm}{\operatorname{lcm}}
\newcommand{\ph}{\varphi}

% The series line printed under every note's title block.
\newcommand{\seriesline}{Erd\H{o}s Problem Notes}

% Production provenance.  The wording is fixed by the corpus provenance
% contract and reproduced verbatim in every manuscript in this repository, so
% the papers cannot drift into different accounts of how they were made.
\newcommand{\noteprovenance}{\provenancenote{\emph{Authorship and AI use.} The
prose, formal proofs, and software in this version were drafted and revised by
large-language-model agents under Will Cook's direction.  Cook set the
objectives and acceptance criteria, reviewed and authorised the manuscript, and
is responsible for its contents.  The agents are production tools, not authors.
See \emph{Statements and declarations}.}}

% The status paragraph every note carries, in its own words, before any result.
% Kernel checking establishes the exact Lean proposition; it does not establish
% that the proposition is the interesting one, that it is new, or that the
% Erdős problem is closed.
\newcommand{\statusboundary}{%
  \emph{Status.}  The problem treated here is open, and this note does not
  close it.  Every statement below marked as checked is a proposition that the
  pinned Lean kernel accepts from the sources this note links to, with no
  \texttt{sorry}, no added axiom, and no unchecked evaluation.  That is a claim
  about the formal statement, not about its mathematical interest, its
  novelty, or the original problem.  The unresolved obligations are named
  exactly, in their own section, and none of the finite computations,
  reductions, or no-go results here removes one of them.}

\hypersetup{
  pdftitle={The prime-gap reformulation and the free carry: Lean-checked finite algebra for Erdős Problem 251},
  pdfsubject={Erdős Problem Notes: Problem 251},
  pdfkeywords={irrationality, prime gaps, dyadic series, summation by parts, Lean 4},
}

\runninghead{Erd\H{o}s \#251: the prime-gap dyadic series}
\title{The Prime-Gap Reformulation and the Free Carry}
\subtitle{Lean-checked finite algebra for Erd\H{o}s Problem~\#251}
\author{Will Cook}
\date{July 2026}

\begin{document}
\maketitle
\noteprovenance

\begin{abstract}
\noindent
Let $p_0=2,p_1=3,\ldots$ enumerate the primes and let
$g_n=p_{n+1}-p_n$.  We formalise in Lean the exact finite identity
\[
 \sum_{i=0}^{n}\frac{p_i}{2^{\,i+1}}
 \;=\;2+\sum_{i=0}^{n-1}\frac{g_i}{2^{\,i+1}}-\frac{p_n}{2^{\,n+1}},
\]
with the endpoint term retained, so that it is available before any statement
about convergence.  For an abstract dyadic tail recurrence
$T_{N+1}=2T_N-g_{N+1}$ with integer digits we prove the block identity
$T_{N+h}=2^{h}T_N-B_{h,N}$, and hence that the shift $T_{N+h}-T_N$ is integral
\emph{if and only if} $(2^{h}-1)T_N$ is; that integrality of one shift
propagates to every later index; and, by Euler's congruence, that a state with
odd reduced denominator $d$ has an integral shift of length $\ph(d)$.

Problem~\#251 is open and nothing here decides it.  Every statement we prove is
finite or algebraic: none of them concerns the infinite series, and connecting
the abstract recurrence to the actual prime-gap tail requires a summability
argument we do not give.  We also record the exact telescoping identity for an
unrestricted integer carry, which is the reason a periodic fractional orbit
carries no information about periodicity of the gap word; the accompanying
construction of prescribed gap patterns is not formalised.  The surviving
obligation is a growing-block anti-concentration statement for actual
consecutive primes.
\end{abstract}

\section{The problem}\label{sec:problem}

Let $p_n$ denote the $n$th prime.  Erd\H{o}s Problem~\#251 asks whether
\[
 \Pi=\sum_{n\ge1}\frac{p_n}{2^{\,n}}
\]
is irrational~\cite{erdos1958}\cite[p.~62]{erdosgraham1980}%
\cite[p.~103]{erdos1988}.  Numbering and status follow Bloom's
catalogue~\cite{erdosproblems}.  The problem is open.

Throughout the formal development the primes are indexed from zero, so that
$p^{(0)}_0=2$, $p^{(0)}_1=3$, $p^{(0)}_2=5$, and the gaps are
$g_i=p^{(0)}_{i+1}-p^{(0)}_i$.  In that indexing
\[
 \Pi=\sum_{i\ge0}\frac{p^{(0)}_i}{2^{\,i+1}} ,
\]
which is the convention every statement below uses; we drop the superscript
from here on.  A second normalisation with denominator $2^{\,i}$ also occurs,
and at every finite horizon it is exactly twice this one,
\[
 \sum_{i=0}^{n-1}\frac{p_i}{2^{\,i}}
 =2\sum_{i=0}^{n-1}\frac{p_i}{2^{\,i+1}} ,
\]
the
\lword{Erdos251/PrimeGapDyadicTail.lean}{55}{prime0DisplayedPartialSumQ_eq_two_mul}{factor-of-two
normalisation}.  A factor of two does not change rationality, so no result
below depends on the choice; the identity is stated so that the indexing cannot
drift silently.

\paragraph{Relation to prior work.}
Erd\H{o}s proved that $\sum_n p_n^{k}/n!$ is irrational for every
$k\ge1$~\cite{erdos1958}, and conjectured both that $\sum_n p_n^{k}/2^{n}$ is
irrational for every $k$ and that $\sum_{n\ge1}p_n/(g_1\cdots g_n)$ is
irrational whenever $g_n\ge2$ and $g_n=o(p_n)$~\cite[p.~103]{erdos1988}; the
choice $g_n=p_n+1$ shows that some growth condition is needed there.  The
statement of Problem~\#251 has been formalised before, as a conjecture with an
unfilled proof, in the \emph{formal-conjectures} collection.  No claim of
priority is made for anything below.

\paragraph{Two routes and where they stop.}
The digits $p_n$ grow, so the series is not a digit expansion in any bounded
alphabet, and the standard rationality criteria for such expansions do not
apply directly.  Two rewritings suggest themselves.  The first replaces the
primes by their gaps, which are bounded on average and are the natural object of
prime-distribution theory; Section~\ref{sec:parts} carries this out exactly, and
identifies the endpoint term that a passage to the limit must control.  The
second reads rationality as a constraint on a tail orbit;
Section~\ref{sec:tail} develops that constraint, and Section~\ref{sec:carry}
gives the exact reason it does not constrain the gaps themselves.

\paragraph{Status of everything below.}
\statusboundary

\begin{table}[t]
\centering
\small
\begin{tabular}{@{}
  >{\raggedright\arraybackslash}p{0.34\linewidth}
  >{\raggedright\arraybackslash}p{0.18\linewidth}
  >{\raggedright\arraybackslash}p{0.40\linewidth}@{}}
\toprule
Statement & Status & Treatment here \\
\midrule
Irrationality of $\Pi$ & Open & Not proved. \\
Finite prime-gap identity & Proved here & Theorem~\ref{res:parts}, with the endpoint retained. \\
Infinite prime-gap identity & Not proved & Requires $p_n2^{-n}\to0$, which we do not establish. \\
Block identity for the tail recurrence & Proved here & Theorem~\ref{res:block}. \\
Integral shift criterion & Proved here; an equivalence & Theorem~\ref{res:shiftiff}. \\
Totient shift from an odd denominator & Proved here & Theorem~\ref{res:totient}. \\
Propagation of an integral shift & Proved here & Theorem~\ref{res:propagate}. \\
Rationality forces periodic gaps & False & Section~\ref{sec:carry}; the carry is free. \\
Anti-concentration for actual gaps & Open & Section~\ref{sec:open}; the surviving obligation. \\
\bottomrule
\end{tabular}
\end{table}

\paragraph{Reading map.}
Section~\ref{sec:parts} is the reformulation, Section~\ref{sec:tail} the tail
algebra, Section~\ref{sec:carry} the no-go, and Section~\ref{sec:open} the
remaining obligation.  Linked phrases open the corresponding Lean declaration
at the pinned source revision~\commitshort.

\medskip
\noindent\textbf{Keywords.} irrationality; prime gaps; dyadic series;
summation by parts; Lean~4.
\qquad
\textbf{MSC 2020.} 11J72 (primary); 11N05, 68V20 (secondary).

\section{Summation by parts, with the endpoint retained}\label{sec:parts}

Write $\Npos$-indexed sums with the convention
\[
 D(P,n)=\sum_{i=0}^{n-1}\frac{P(i)}{2^{\,i+1}},
 \qquad
 \Delta(P,n)=\sum_{i=0}^{n-1}\frac{P(i+1)-P(i)}{2^{\,i+1}} ,
\]
for a rational sequence $P$.  These are the
\lword{Erdos251/PrimeGapDyadicTail.lean}{45}{dyadicPartialSumQ}{dyadic partial
sum} and the
\lword{Erdos251/PrimeGapDyadicTail.lean}{65}{dyadicDifferencePartialSumQ}{dyadic
difference sum}.

\begin{proposition}[finite summation by parts]\label{res:abel}
For every rational sequence $P$ and every $n\ge0$,
\[
 D(P,n+1)=P(0)+\Delta(P,n)-\frac{P(n)}{2^{\,n+1}} .
\]
\end{proposition}

\begin{proof}
Induction on $n$.  At $n=0$ both sides equal $P(0)/2$.  For the step, adding
$P(n+1)/2^{\,n+2}$ to the left and
$\bigl(P(n+1)-P(n)\bigr)/2^{\,n+1}$ to the difference sum changes the endpoint
term from $P(n)/2^{\,n+1}$ to $P(n+1)/2^{\,n+2}$, and the two adjustments agree.
\end{proof}

\noindent Formalised as the
\lword{Erdos251/PrimeGapDyadicTail.lean}{82}{dyadicPartialSumQ_eq_start_add_differences}{summation-by-parts
identity}.  Nothing is assumed about $P$: no positivity, no monotonicity, and no
convergence.

Specialising to $P(i)=p_i$, whose first value is $p_0=2$, and writing
$g_i=p_{i+1}-p_i$ for the zero-based gaps, gives the reformulation.

\begin{theorem}[prime-gap reformulation]\label{res:parts}
For every $n\ge0$,
\[
 \sum_{i=0}^{n}\frac{p_i}{2^{\,i+1}}
 =2+\sum_{i=0}^{n-1}\frac{g_i}{2^{\,i+1}}-\frac{p_n}{2^{\,n+1}} .
\]
\end{theorem}

\noindent Formalised as the
\lword{Erdos251/PrimeGapDyadicTail.lean}{111}{prime0_dyadic_summation_by_parts}{prime-gap
summation by parts}, using the
\lword{Erdos251/PrimeGapDyadicTail.lean}{105}{primeGapPartialSumQ}{gap partial
sum} and the
\lword{Erdos251/PrimeGapDyadicTail.lean}{100}{primeGap0_cast}{gap cast
identity}; the latter records that the natural-number difference
$p_{n+1}-p_n$ agrees with the difference taken in $\Q$, which needs
$p_n\le p_{n+1}$, the
\lword{Erdos251/PrimeGapDyadicTail.lean}{96}{prime0_mono_step}{monotonicity
step}.  The leading $2$ is the first prime, not a normalising constant.

\paragraph{What the identity does and does not give.}
The identity is exact at every finite horizon and free of hypotheses.  It is
not the corresponding statement about infinite series.  Passing to the limit
requires $p_n2^{-n}\to0$, and while that is elementary, it is not proved in the
formal source and is therefore not available to any statement here.  We record
the finite identity in this form deliberately: the endpoint term is the whole
content of the passage to the limit, and suppressing it is exactly the step
that would make a finite identity look like an infinite one.

\section{The tail recurrence and integral shifts}\label{sec:tail}

Rationality of a dyadic series constrains its tail orbit.  This section proves
what that constraint is, for an abstract recurrence.

\begin{definition}\label{def:rec}
Let $g:\N\to\Z$ and $T:\N\to\Q$.  Say $T$ satisfies the \emph{dyadic tail
recurrence} with digits $g$ when
\[
 T_{N+1}=2T_N-g_{N+1}\qquad\text{for every }N .
\]
Write $\sigma_h(N)=T_{N+h}-T_N$ for the \emph{shift} of length $h$ at $N$, and
call a rational number \emph{integral} when it is the image of an integer.
\end{definition}

\noindent These are the
\lword{Erdos251/PrimeGapDyadicTail.lean}{123}{DyadicTailRecurrence}{tail
recurrence}, the
\lword{Erdos251/PrimeGapDyadicTail.lean}{127}{tailShift}{shift}, and
\lword{Erdos251/PrimeGapDyadicTail.lean}{143}{RatIntegral}{integrality}.
The digits are arbitrary integers.  In the intended instance they are the prime
gaps and $T_N$ is the scaled tail of $\Pi$ after level $N$, but that
instantiation needs a summability argument and is not made here; every
statement below is a theorem about Definition~\ref{def:rec}.

Iterating the recurrence accumulates an explicit integer.  Define
$B_{0,N}=0$ and $B_{h+1,N}=2B_{h,N}+g_{N+h+1}$, so that
$B_{h,N}=g_{N+1}2^{\,h-1}+\cdots+g_{N+h}$: the
\lword{Erdos251/PrimeGapDyadicTail.lean}{149}{dyadicTailBlock}{tail block}.

\begin{theorem}[block identity]\label{res:block}
For every $N$ and $h$,
\[
 T_{N+h}=2^{h}T_N-B_{h,N},
 \qquad\text{hence}\qquad
 \sigma_h(N)=(2^{h}-1)\,T_N-B_{h,N} .
\]
\end{theorem}

\begin{proof}
Induction on $h$; the step is one application of the recurrence together with
$2\cdot2^{h}=2^{h+1}$ and the recursion defining $B$.  The second identity is
the first minus $T_N$.
\end{proof}

\noindent Formalised as the
\lword{Erdos251/PrimeGapDyadicTail.lean}{155}{tail_iterate_eq_pow_mul_sub_block}{iterated
block identity} and the
\lword{Erdos251/PrimeGapDyadicTail.lean}{169}{tailShift_eq_scaled_sub_block}{scaled
shift identity}.  The shift also obeys the recurrence in its own right,
$\sigma_h(N+1)=2\sigma_h(N)-(g_{N+h+1}-g_{N+1})$, the
\lword{Erdos251/PrimeGapDyadicTail.lean}{131}{tailShift_succ}{shift step
identity}.

Since $B_{h,N}$ is an integer, Theorem~\ref{res:block} converts a question about
the shift into a question about a single scaled state.

\begin{theorem}[integral-shift criterion]\label{res:shiftiff}
For every $N$ and $h$, the shift $\sigma_h(N)$ is integral if and only if
$(2^{h}-1)T_N$ is integral.
\end{theorem}

\begin{proof}
By Theorem~\ref{res:block} the two differ by the integer $B_{h,N}$, and
subtracting an integer does not change integrality.
\end{proof}

\noindent Formalised as the
\lword{Erdos251/PrimeGapDyadicTail.lean}{225}{tailShift_integral_iff_scaledTail}{integral-shift
criterion}, on the
\lword{Erdos251/PrimeGapDyadicTail.lean}{179}{ratIntegral_sub_int_iff}{integer-shift
invariance}.  This is an equivalence, not a one-way reduction: the block
carries no information about integrality, so the two conditions are the same
condition.

\begin{theorem}[a shift of totient length]\label{res:totient}
If the reduced denominator $d$ of $T_N$ is odd, then $\sigma_{\ph(d)}(N)$ is
integral.
\end{theorem}

\begin{proof}
Since $d$ is odd, $2$ and $d$ are coprime, so Euler's congruence gives
$2^{\ph(d)}\equiv1\pmod d$, that is $d\mid 2^{\ph(d)}-1$.  Writing
$2^{\ph(d)}-1=dk$ and $T_N=u/d$ in lowest terms, $(2^{\ph(d)}-1)T_N=ku$ is an
integer, and Theorem~\ref{res:shiftiff} transfers this to the shift.
\end{proof}

\noindent Formalised as the
\lword{Erdos251/PrimeGapDyadicTail.lean}{236}{tailShift_integral_totient_of_odd_den}{totient
shift} and the
\lword{Erdos251/PrimeGapDyadicTail.lean}{197}{ratIntegral_totientMultiplier_of_odd_den}{Euler
multiplier}.  The hypothesis is a genuine restriction: the argument uses
coprimality of $2$ with the denominator, and the even part of a denominator is
exactly what the doubling in the recurrence acts on.

\begin{theorem}[propagation]\label{res:propagate}
If $\sigma_h(N)$ is integral, then $\sigma_h(N+k)$ is integral for every
$k\ge0$.
\end{theorem}

\begin{proof}
By the shift step identity, $\sigma_h(N+1)=2\sigma_h(N)-(g_{N+h+1}-g_{N+1})$ is
an integer combination of an integer and two digits.  Induct on $k$.
\end{proof}

\noindent Formalised as the
\lword{Erdos251/PrimeGapDyadicTail.lean}{247}{tailShift_integral_succ}{one-step
propagation} and the
\lword{Erdos251/PrimeGapDyadicTail.lean}{259}{tailShift_integral_add}{propagation
to every later index}.

\paragraph{What this section proves, and what it does not.}
Theorems~\ref{res:block}--\ref{res:propagate} are statements about
Definition~\ref{def:rec} and hold for arbitrary integer digits.  None of them
mentions primes, and none of them uses convergence.  Instantiating them at the
prime-gap tail is a separate step: one must first prove that the scaled tail of
$\Pi$ exists and satisfies the recurrence, and we do not.  In particular
nothing here shows that $\Pi$ has, or does not have, an integral shift.

\section{The carry is free}\label{sec:carry}

A natural hope is that rationality of a dyadic series forces its digit stream
to be eventually periodic, or at least automatic, and that this is impossible
for prime gaps.  The hope fails at the first step, and the reason is a
one-line identity.

Let $K:\N\to\Q$ be arbitrary and put
$\kappa_n=2K_n-K_{n+1}$: the
\lword{Erdos251/PrimeGapDyadicTail.lean}{270}{carryCoeff}{carry coefficient}.

\begin{proposition}[exact telescoping]\label{res:telescope}
For every $n\ge0$,
\[
 \sum_{i=0}^{n-1}\frac{\kappa_i}{2^{\,i+1}}=K_0-\frac{K_n}{2^{\,n}} .
\]
\end{proposition}

\begin{proof}
Induction on $n$; the added term is
$(2K_n-K_{n+1})/2^{\,n+1}=K_n/2^{\,n}-K_{n+1}/2^{\,n+1}$.
\end{proof}

\noindent Formalised as the
\lword{Erdos251/PrimeGapDyadicTail.lean}{279}{carryPartialSum_eq}{carry
telescoping identity}, on the
\lword{Erdos251/PrimeGapDyadicTail.lean}{274}{carryPartialSum}{carry partial
sum}.  When the carries are natural numbers and $K_{n+1}\le2K_n$, the
coefficients are natural numbers and the two readings agree, the
\lword{Erdos251/PrimeGapDyadicTail.lean}{295}{natCarryCoeff_cast}{natural-carry
cast}.

\paragraph{The reading, and its exact status.}
The identity says that the carry sequence $K$ is unconstrained: choose $K$
freely, and the emitted coefficients $\kappa_n$ have dyadic partial sums
telescoping to $K_0$ up to $K_n2^{-n}$.  A rational limit therefore places no
periodicity or automaticity requirement on the coefficient stream, because the
same rational value arises from carry sequences that are not eventually
periodic.  This is why the tail constraint of Section~\ref{sec:tail} does not
descend to the gap word: rationality controls a fractional orbit, and the
integer part absorbs everything else.

Proposition~\ref{res:telescope} is proved; the accompanying claim that
prescribed finite gap patterns can be realised by such a carry is a
construction we have not formalised, and it is not used anywhere above.  What
is established is the identity and the consequence that the periodicity route
is closed, not a classification of which coefficient streams occur.

\section{Complements and further questions}\label{sec:open}

Problem~\#251 is open.  The results above fix the shape of what is missing.

\begin{problem}[growing-block anti-concentration]
Show that for the actual sequence of consecutive prime gaps, every
sufficiently long block stays a definite dyadic distance from every
carry-compatible rational target, with the required length growing with the
starting index.
\end{problem}

\begin{problem}[the input from prime distribution]
Identify a theorem about consecutive primes strong enough to supply that block
obstruction.  The occurrence of individual gap values, however often, is not
enough: by Section~\ref{sec:carry} a finite insertion of any fixed even gap
word is compatible with a rational value, so the obstruction must be a
statement about long blocks rather than about single gaps.
\end{problem}

Two smaller obligations remain on the formal side.  The infinite form of
Theorem~\ref{res:parts} needs $p_n2^{-n}\to0$; and instantiating
Definition~\ref{def:rec} at the prime-gap tail needs summability.  Neither is
difficult, and neither is done, so no statement above may be read as concerning
the infinite series.  What the formal content changes is the location of the
difficulty: it is not in the reformulation, and it is not in the tail algebra.

\pdfbookmark[1]{Statements and declarations}{statements-declarations}
\subsection*{Statements and declarations}

\paragraph{Artefact and data availability.} The
\href{\sourceurl}{pinned formal-source revision} contains the Lean sources, the
fixed toolchain, and the library manifest used in the verification.  This
manuscript provides navigation rather than proof authority.

\paragraph{Declaration of generative AI use.} Large-language-model agents were
used throughout development to draft and revise prose, formal proofs, and
software.  The author set the objectives and acceptance criteria, selected and
reviewed the claims, and approved the published version.  The author assumes
responsibility for the accuracy, interpretation, and presentation of the work.
Generative systems are production tools; they are not authors and supply no
independent authority.  Lean checks each proof term against the fixed library
version, and the sources linked here contain no proof placeholders and no
project-defined axioms; Lean does not authorise the exposition, the citation
choices, or the interpretation, for which the author remains responsible.

\paragraph{Funding and competing interests.} This work received no external
funding.  The author declares no competing interests.

\paragraph{Acknowledgements.} The problem numbering and status follow the
Erd\H{o}s Problems catalogue maintained by Thomas Bloom~\cite{erdosproblems}.

\appendix
\section{Guide to the formal sources}\label{app:index}

Each linked phrase opens its Lean declaration at the pinned source
revision~\commitshort.  All declarations of this note live in one module.  The
prime-specific material is confined to Section~\ref{sec:parts}; everything in
Sections~\ref{sec:tail} and~\ref{sec:carry} is stated for arbitrary integer
digits and arbitrary rational carries, and the reader should not transfer those
statements to the primes without the summability step named in
Section~\ref{sec:open}.

% The exhaustive source manifest is retained in the authored TeX for automated
% validation, and kept outside the printed note.
\iffalse

\begin{tabular}{@{}ll@{}}
informal statement & linked Lean source\\
zero-based primes & \lrefx{Erdos251/PrimeGapDyadicTail.lean}{29}{prime0}\\
zero-based gaps & \lrefx{Erdos251/PrimeGapDyadicTail.lean}{33}{primeGap0}\\
first gap value & \lrefx{Erdos251/PrimeGapDyadicTail.lean}{36}{primeGap0_zero}\\
second gap value & \lrefx{Erdos251/PrimeGapDyadicTail.lean}{40}{primeGap0_one}\\
dyadic partial sum & \lrefx{Erdos251/PrimeGapDyadicTail.lean}{45}{dyadicPartialSumQ}\\
displayed normalisation & \lrefx{Erdos251/PrimeGapDyadicTail.lean}{50}{prime0DisplayedPartialSumQ}\\
factor-of-two identity & \lrefx{Erdos251/PrimeGapDyadicTail.lean}{55}{prime0DisplayedPartialSumQ_eq_two_mul}\\
dyadic difference sum & \lrefx{Erdos251/PrimeGapDyadicTail.lean}{65}{dyadicDifferencePartialSumQ}\\
partial-sum step & \lrefx{Erdos251/PrimeGapDyadicTail.lean}{68}{dyadicPartialSumQ_succ}\\
difference-sum step & \lrefx{Erdos251/PrimeGapDyadicTail.lean}{73}{dyadicDifferencePartialSumQ_succ}\\
summation by parts & \lrefx{Erdos251/PrimeGapDyadicTail.lean}{82}{dyadicPartialSumQ_eq_start_add_differences}\\
monotonicity step & \lrefx{Erdos251/PrimeGapDyadicTail.lean}{96}{prime0_mono_step}\\
gap cast identity & \lrefx{Erdos251/PrimeGapDyadicTail.lean}{100}{primeGap0_cast}\\
gap partial sum & \lrefx{Erdos251/PrimeGapDyadicTail.lean}{105}{primeGapPartialSumQ}\\
prime-gap summation by parts & \lrefx{Erdos251/PrimeGapDyadicTail.lean}{111}{prime0_dyadic_summation_by_parts}\\
tail recurrence & \lrefx{Erdos251/PrimeGapDyadicTail.lean}{123}{DyadicTailRecurrence}\\
shift & \lrefx{Erdos251/PrimeGapDyadicTail.lean}{127}{tailShift}\\
shift step identity & \lrefx{Erdos251/PrimeGapDyadicTail.lean}{131}{tailShift_succ}\\
integrality & \lrefx{Erdos251/PrimeGapDyadicTail.lean}{143}{RatIntegral}\\
tail block & \lrefx{Erdos251/PrimeGapDyadicTail.lean}{149}{dyadicTailBlock}\\
iterated block identity & \lrefx{Erdos251/PrimeGapDyadicTail.lean}{155}{tail_iterate_eq_pow_mul_sub_block}\\
scaled shift identity & \lrefx{Erdos251/PrimeGapDyadicTail.lean}{169}{tailShift_eq_scaled_sub_block}\\
integer-shift invariance & \lrefx{Erdos251/PrimeGapDyadicTail.lean}{179}{ratIntegral_sub_int_iff}\\
Euler multiplier & \lrefx{Erdos251/PrimeGapDyadicTail.lean}{197}{ratIntegral_totientMultiplier_of_odd_den}\\
integral-shift criterion & \lrefx{Erdos251/PrimeGapDyadicTail.lean}{225}{tailShift_integral_iff_scaledTail}\\
totient shift & \lrefx{Erdos251/PrimeGapDyadicTail.lean}{236}{tailShift_integral_totient_of_odd_den}\\
one-step propagation & \lrefx{Erdos251/PrimeGapDyadicTail.lean}{247}{tailShift_integral_succ}\\
propagation to every later index & \lrefx{Erdos251/PrimeGapDyadicTail.lean}{259}{tailShift_integral_add}\\
carry coefficient & \lrefx{Erdos251/PrimeGapDyadicTail.lean}{270}{carryCoeff}\\
carry partial sum & \lrefx{Erdos251/PrimeGapDyadicTail.lean}{274}{carryPartialSum}\\
carry telescoping identity & \lrefx{Erdos251/PrimeGapDyadicTail.lean}{279}{carryPartialSum_eq}\\
natural carry coefficient & \lrefx{Erdos251/PrimeGapDyadicTail.lean}{292}{natCarryCoeff}\\
natural-carry cast & \lrefx{Erdos251/PrimeGapDyadicTail.lean}{295}{natCarryCoeff_cast}\\
\end{tabular}
\fi

\begin{thebibliography}{9}
\small
\raggedright
\setlength{\itemsep}{0pt}
\bibitem{erdos1958} P.~Erd\H{o}s,
  \href{https://users.renyi.hu/~p_erdos/1958-08.pdf}{\emph{Sur certaines
  s\'eries \`a valeur irrationnelle}},
  Enseign. Math. (2) \textbf{4} (1958), 93--100. MR~98732.
\bibitem{erdosgraham1980} P.~Erd\H{o}s and R.~L. Graham,
  \href{https://mathweb.ucsd.edu/~ronspubs/80_11_number_theory.pdf}{\emph{Old
  and New Problems and Results in Combinatorial Number Theory}},
  Monogr. Enseign. Math. 28, Geneva, 1980, p.~62.
\bibitem{erdos1988} P.~Erd\H{o}s, \emph{On the irrationality of certain series:
  problems and results}, in A.~Baker (ed.), \emph{New Advances in Transcendence
  Theory}, Cambridge UP, 1988, pp.~102--109,
  doi:\href{https://doi.org/10.1017/CBO9780511897184.009}{10.1017/CBO9780511897184.009}.
\bibitem{erdosproblems} T.~F. Bloom,
  \href{https://www.erdosproblems.com/251}{\emph{Erd\H{o}s Problem \#251}},
  \texttt{erdosproblems.com/251}, accessed 22 July 2026.
\end{thebibliography}

\end{document}
